\section{Authorship and use of AI}\label{app:contrib}

\paragraph{Course project.} The first version of this study was the ME~306 course project of a team of five: Siddhant Patil (24B2121), Vaishnavi Shivai (24B2116), Avinav Shekhar (24B2186), Anshuman Gupta (24B2169) and Vidit Ganna (24B2210). The suggested division into separate roles was not followed strictly; the learning goal was for every member to understand the whole study, so a single common method was worked out together and each part was discussed jointly. That version is preserved as tag \path{rebuild} of the repository (Appendix~\ref{app:code}).

\paragraph{This version.} After the submission, Siddhant Patil extended the study: the automatic refinement to convergence, the strict ranking under the literal reading of the \SI{7.7}{MPa} limit, the nine studies of Section~\ref{sec:studies}, the treatment of assumptions in Section~\ref{sec:limits}, and this report. ``We'' in the text refers to the study as a whole.

\paragraph{Use of AI.} Claude (Anthropic) was used as an assistant in three ways: writing and running the code under direction, collecting and checking references, and helping to prepare this report. The code was built one step at a time, and each step was reviewed and tested before the next (Appendix~\ref{app:validation}).

\paragraph{Decisions.} All modelling decisions were taken by the authors. These include how the given limits are read, how the recuperators and coolers are specified, how the baseline mass is calibrated, which objectives are used, how the final design is selected and how each study is set up (Appendix~\ref{app:decisions}). Each was reached through discussion, including arguing against the AI's suggestions and having it argue against ours, until the reasoning and its caveats were understood. Several of its first suggestions were rejected or changed in this way. The results were also checked independently of the main solver (Appendix~\ref{app:validation}). Responsibility for the method, the results and any remaining errors lies with the authors.
